\documentclass[final,1p,times]{elsarticle}
\usepackage{amsmath,amssymb,amsthm,mathtools}
\usepackage{booktabs,multirow}
\usepackage{graphicx}
\usepackage[colorlinks=true,linkcolor=blue,citecolor=blue]{hyperref}
\usepackage{placeins}
% Relax float placement so figures/tables are not deferred to the end
% of the document (where they would land inside the references).
\renewcommand{\topfraction}{0.85}
\renewcommand{\textfraction}{0.1}
\renewcommand{\floatpagefraction}{0.85}
\setcounter{topnumber}{4}
\setcounter{totalnumber}{6}

\newtheorem{theorem}{Theorem}
\newtheorem{lemma}{Lemma}
\newtheorem{proposition}{Proposition}
\newtheorem{definition}{Definition}
\newtheorem{remark}{Remark}
\newcommand{\etal}{\emph{et al.}}
\newcommand{\opt}{\mathrm{opt}}

\begin{document}
\begin{frontmatter}
\title{Supplementary Material for\\
``Noise, not optimality: a capacity-to-budget diagnostic for
stealthy poisoning in differentially private federated learning''}
\end{frontmatter}

\section{Proofs}\label{app:proofs}
%======================================================================
\textbf{Setting.} In the isotropic regime ($h{=}0$) the benign update
is $x_b=m+\xi$ with $\xi\sim\mathcal{N}(0,\sigma^2 I_d)$, and the
attacker displaces one update by $\Delta$, $\|\Delta\|=\rho$. The
server observes the statistic $s(x)=\|x-\hat m\|^2$, where $\hat m$ is
the coordinate median of the benign population; for symmetric
dispersion, $\hat m$ is consistent for $m$, so we condition on
$\hat m=m$.

\textbf{(i) Attacker statistic.} Writing
$s(x_b+\Delta)=\sum_{j=1}^{d}(\xi_j+\Delta_j)^2$, each
$(\xi_j+\Delta_j)^2/\sigma^2$ is a noncentral $\chi^2_1$ with
noncentrality $\lambda_j=\Delta_j^2/\sigma^2$. Hence
\[
  \mathbb{E}[s]=\sigma^2\textstyle\sum_j(1+\lambda_j)=d\sigma^2+\rho^2,
\]
and $\mathrm{Var}[s]=\sigma^4\sum_j 2(1+2\lambda_j)=2\sigma^4
d+4\sigma^2\rho^2$. In the stealthy regime $\rho^2\le d\sigma^2$ the
second term is $O(1/d)$ relative to the first, so
$\mathrm{std}[s]=\sigma^2\sqrt{2d}\,(1+o(1))$, a dominant-term
(delta-method) argument; this is where the $\sqrt{d}$ dispersion in
the denominator originates.

\textbf{(ii) Benign statistic.} The same computation with $\Delta=0$
gives $\mathbb{E}[s_b]=d\sigma^2$ and the same standard deviation to
leading order. Therefore
\[
  \mathrm{sep}
  =\frac{\mathbb{E}[s]-\mathbb{E}[s_b]}{\mathrm{std}[s]}
  =\frac{\rho^2}{\sigma^2\sqrt{2d}}
  =c\,\frac{\rho^2}{\sigma^2\sqrt{d}},
  \qquad c=\tfrac{1}{\sqrt{2}},
\]
which is the separation formula with $w^2=\sigma^2$, and inverting
$\mathrm{sep}=\mathrm{sep}^{*}$ yields the stealth-capacity bound with
$\eta_{\max}=\sqrt{\mathrm{sep}^{*}/c}\;\sigma\,d^{1/4}$.

\textbf{(iii) From surrogate to pipeline.} With a benign
heterogeneity floor, Lemma~1 replaces $\sigma^2$ by
$w^2=\sigma^2+h^2$, giving the two-regime law of
Proposition~1; the isotropic slice $h=0$ recovers
the theorem. The measured collapse constant $c=0.30$ is smaller than
the surrogate value $1/\sqrt{2}$ because the median of $N$ finite
benign samples absorbs part of the shift and the empirical threshold
carries finite-sample dispersion (Remark~1); the
functional form is unaffected, which is why the defender optimization
calibrates the constant rather than deriving it. \hfill$\square$

\textbf{Proof of Theorem~2.} The Lagrangian of
the budget constraint separates across clients; stationarity gives
$p_i=\min\{\eta_{\max,i},\alpha_i/(2\mu)\}$ \cite{boyd2004convex}, where
$\mu\ge0$ is the Lagrange multiplier on the budget constraint; the
constraint is \emph{slack} ($\mu{=}0$) exactly when $C<\Gamma$ (case~1, all
caps saturated, budget slack). For case~2 the exchange argument holds
only under the stated cap-saturation condition: moving a unit of
budget to an unselected client $i$ improves the objective only if
$i$'s first increment is uncapped (marginal value $\alpha_i)$
\emph{and} exceeds the marginal value of the capped increment on the
donor; when the donor is itself capped its marginal value is
$\alpha_j$, so the rule favours the higher-$\alpha$ capped client, not
the higher-$\alpha\eta_{\max}$ one. Bisection on $\mu$ (or the
$\mu{=}0$ trivial branch) is the practical solver. \hfill$\square$

\textbf{Proof of Theorem~3.} $D$ is non-decreasing in
each $p_i$; for fixed $t$ the constraint $g\ge t$ is a monotone
covering constraint, satisfied at minimum L2 cost by
the covering-cost expression; maximizing $t$ subject to $B(t)\le\Gamma$ yields
$t^\star$. \hfill$\square$

\textbf{Proof of Theorem~4.} The attack is deployable iff
$K\,\eta_{\max}^2(\sigma)\ge\Gamma$. Substituting the stealth-capacity bound
and solving for $\sigma$ gives
$\sigma\ge\sqrt{c\,\Gamma/(\mathrm{sep}^{*}K)}\;d^{-1/4}$, i.e.\
$\sigma_{\max}=\Theta(\sqrt{\Gamma/K}\,d^{-1/4})$; since
$\eta_{\max}$ is strictly increasing in $\sigma$, the stealthy-escape
region grows monotonically with the privacy level. \hfill$\square$

%======================================================================
%======================================================================
\section*{Supplementary Material}
%======================================================================
The following material supports the main text but is not required to
follow the argument; it is provided for reproducibility and to document
robustness checks.

\subsection{Proxy framework: capacity law and allocator
ablation}\label{supp:proxy}
This section contains the proxy-level evaluation of the main text
(Experimental Evaluation): setup, mechanism-faithful baselines, full
$\sigma$/$\Gamma$ sweeps, and the allocator ablation.

\subsubsection{Setup}
Coordinate median; $d{=}1500$, $d_0{=}100$, $N{=}30$, $K{=}8$;
heterogeneous leverage $\alpha_i\sim\mathrm{LogN}(0,0.8)$; capacities
$\eta_{\max,i}=14.5\,\sigma\, u_i$, $u_i\sim U[0.5,1.5]$ (this is the
isotropic-statistics regime of Theorem~1, \emph{instantiated}
so that the allocator can be exercised under the analytically derived
law). All baselines share
the attacker's budget envelope $\sum p_i^2\le\Gamma$ and differ only in
their information assumptions.

\subsubsection{Baselines (mechanism-faithful proxies)}
\begin{itemize}
  \item \textbf{S21} \cite{shejwalkar2021turning}: DP-unaware; solves
  water-filling with capacity misjudged at $\sigma_{\mathrm{ref}}{=}1$.
  \item \textbf{W25} \cite{wang2025ldpfl}: noise-adaptive global
  scaling $p_i=\min\{\kappa\sigma,\sqrt{\Gamma/K}\}$, client-blind---a
  static reduction of Wang \etal's LDP-aware attack in which the
  tail-aware direction refinement is replaced by the common direction
  $-\hat g$ shared by all allocators, so that direction effects are
  isolated and compared separately in the robustness section below.
  \item \textbf{H25} \cite{hossain2025mpelm}: profile clipping
  $p_i=\min\{\eta_i,\lambda\alpha_i\}$ without redirection.
  \item \textbf{Q25} \cite{qiu2025budget}: uncapped
  $p_i=\lambda\alpha_i$ (its global-budget allocation principle,
  ported from non-FL data poisoning to our setting).
\end{itemize}

\subsubsection{Results (proxy capacity analysis)}
The full $\sigma$ and $\Gamma$ sweeps over the six allocators
(repository CSV; ablation in Fig.~\ref{fig:6}) underlie the findings
below: each baseline collapses in a distinct regime; ours never does.

\textbf{Findings (proxy).}
(1) Distinct failure modes: S21 overestimates capacity at
$\sigma<\sigma_{\mathrm{ref}}$ (total collapse at $\sigma{=}0.5$) and
under-spends at $\sigma>1$ (plateau at $271$, $-44\%$ vs.\ ours at
$\sigma{=}8$); Q25 collapses ($-67\%$ at $\Gamma{=}16000$);
W25 under-spends at low $\sigma$; H25 loses $18\%$ at
$\Gamma{=}16000$ (no redirection).
(2) In the uncapped regime ($\sigma\ge4$, moderate $\Gamma$) Q25
coincides with water-filling within $0.7\%$; the residual gap is
the selection criterion ($\alpha$ vs.\ $\alpha\eta$), consistent with
the ablation below.
(3) Ablation (Fig.~\ref{fig:6}): removing capping costs $66.4\%$ on average,
removing client selection $53.1\%$; at small $\Gamma$ pure-$\alpha$
selection is (slightly) optimal; the $\alpha\eta$ criterion pays
off when capping binds.
(4) Scale robustness: re-running the full $\sigma$ sweep at
$(N,K){=}(20,6),(30,8),(50,12),(100,20)$ ($10$ scenario seeds per cell,
$\Gamma{=}500K$ so the per-client budget is held fixed) leaves every
qualitative conclusion unchanged -- the collapse pattern is identical at
all four scales (S21 and W25 fail in all six $\sigma$ cells; H25, Q25,
and UNI in four), and OURS never falls below the best baseline by more
than $3\%$ in any cell (the residual gap is the known uncapped-regime
equivalence of Q25). The full grid is in the repository CSV.

\begin{table}[!t]
\centering\small
\caption{Hyperparameters and protocol of the real FL pipeline (shared
across FashionMNIST, MNIST, and CIFAR-10 unless noted).}
\label{tab:hyper}
\begin{tabular}{@{}p{0.28\linewidth}p{0.64\linewidth}@{}}
\hline
Model (FMNIST/MNIST) & $16{\times}3^2$ conv--maxpool--
$32{\times}3^2$ conv--maxpool--FC; $d{=}20{,}490$ \\
Model (CIFAR-10) & $32{\times}3^2$ conv--maxpool--
$64{\times}3^2$ conv--maxpool--FC; $d{=}60{,}362$ \\
Local training & SGD, lr $0.08$ (no momentum), batch $64$; $2$ local
epochs (warm-up), $1$ local epoch (attack rounds) \\
Warm-up & $60$ rounds, size-weighted FedAvg; $2{,}000$ validation
samples; test acc.\ $0.846$ (FMNIST), $0.977$ (MNIST), $0.607$
(CIFAR-10) \\
Federation & $N{=}20$ clients of $300$--$900$ samples each; $K{=}6$
compromised \\
Noise scale & $\sigma_{\mathrm{ben}}$ = mean over coordinates of the
per-coordinate std.\ of the benign update pool;
$\sigma_{\mathrm{rel}}{=}\sigma/\sigma_{\mathrm{ben}}$ \\
Detection & per-round threshold = mean ${+}\,1.645$ std.\ of the benign
pool ($\approx 5\%$ FPR); flagged clients removed \\
Grids & R1: $\sigma_{\mathrm{rel}}\in\{0.5,1,2,4,8,16,24,32,48,64\}$; R2/R2b:
$\sigma_{\mathrm{rel}}\in\{0.5,1,2\}\times
\kappa\in\{0.25,0.5,1,2\}$, $\Gamma{=}\kappa K\|\Delta\|^2$ \\
Repetition & $5$ attack rounds $\times$ $10$ seeds per cell;
$\eta_{\max,i}$ by bisection ($6$ trials) on the deployment detector \\
\hline
\end{tabular}
\end{table}

\begin{figure}[!t]
\centering
\includegraphics[width=0.8\linewidth]{figures/fig6_ablation.pdf}
\caption{Ablation on the proxy framework: hidden impact vs.\ budget
$\Gamma$ for full water-filling and its degraded variants. Removing the
capping term costs $66.4\%$ on average; removing client selection
$53.1\%$; the budget-relaxation variant tracks full except at mid-range
$\Gamma$.}
\label{fig:6}
\end{figure}


\subsection{Central-DP calibration and the $\varepsilon(\sigma)$ mapping}\label{supp:eps}
For the Gaussian mechanism with L2 sensitivity $\Delta$ (the full
per-update norm) and per-coordinate noise scale $\sigma$, the analytic
Gaussian mechanism gives $\varepsilon = 2\Delta^2\ln(1.25/\delta)/\sigma^2$
at a fixed $\delta$~\cite{dwork2006calibrating}. In our pipeline
$\Delta\approx0.36$ and $\sigma=\sigma_{\mathrm{rel}}\cdot0.0019$, so at
the studied $\sigma_{\mathrm{rel}}\in\{0.5,\dots,16\}$ the per-round
$\varepsilon\gg1$: the noise is calibrated to the \emph{per-coordinate}
detection scale rather than to a meaningful per-update DP guarantee, and
we therefore demonstrate the detection-masking \emph{mechanism} at sub-DP
noise scales (a defender requiring $\varepsilon\ll1$ would need $\sigma$
deep in the noise-dominated regime, compounding the utility loss). We
report $\varepsilon(\sigma)$ alongside $\sigma$ with the real-pipeline
capacity probes. Under \emph{true} central-DP-on-the-aggregate (noise
added only after aggregation) the server could run detection on the clean
pre-noise aggregate and the masking term would vanish; all results are
therefore statements about the per-update noise model. That regime is not
a strawman: robust aggregators are nonlinear, so DP-on-the-aggregate does
not compose with them---the coordinate median is not a low-sensitivity
query, and noise added after it is not a calibrated DP
mechanism---whereas deployments that stack DP with robust aggregation
noise client updates before the aggregator, as in client-level DP
training~\cite{mcmahan2018dp} and per-client auditing or adaptive-clipping
pipelines; the main text flags the boundary explicitly wherever a
conclusion depends on it.

\subsection{Robustness: attack direction, aggregation rule, adaptive
attacker, and a third dataset (R2c--R5)}\label{sec:robust}
\subsubsection{R2c: robustness to attack-direction sophistication (SOTA attacks)}
\label{sec:r2c}
All allocators above inject along a single default direction, the
negative validation gradient $\hat g$ (gradient-ascent on loss), the
standard model-poisoning direction. We therefore test robustness against a
\emph{stronger, aggregator-tailored} attack direction, closing with a
head-to-head against the state-of-the-art optimal/tailored attacks.

We instantiate, inside the same pipeline, the attack directions of
Shejwalkar and Houmansadr (NDSS 2021)~\cite{shejwalkar2021turning}, which are
also the mechanism-level embodiments of the two adjacent works we compare
against:
\begin{itemize}
  \item \texttt{shej\_sign} --- coordinate sign of the benign aggregate
  (the sign-flip direction of Shejwalkar's AGR-tailored attack, which
  targets median-type aggregators; it is also the closest static
  mechanism-level proxy for the tail-aware adaptive directions of
  Wang et al.'s LDPFL-Attack~\cite{wang2025ldpfl});
  \item \texttt{shej\_std} --- benign coordinate variance weighting
  (Shejwalkar's variance-tailored direction);
  \item \texttt{shej\_unit} --- the benign reference unit vector
  (Shejwalkar's AGR-agnostic direction);
  \item \texttt{grad} --- our default $-\hat g$.
\end{itemize}
The three Shejwalkar directions are extracted from the authors' official
NDSS 2021 implementation (\texttt{our\_attack\_median}).
Each direction is held to the same DP-masked stealth cap
$\eta_{\max,i}$ (recomputed per direction with the same probing
protocol and the same attacker information requirements as the default
direction, so switching direction changes no information assumption;
Hossain's rule of hiding the
attack within the DP-noise envelope~\cite{hossain2025mpelm} is the same cap, already used by
all allocators). We re-run OURS (water-filling) and Q25 (uncapped budget,
the frontier occupant on CIFAR-10) under all four directions across the
same $\sigma\in\{0.5,1,2\}$, $\kappa\in\{0.5,2\}$, $10$ seeds.

\emph{Finding (direction does not rescue stealth-aware optimality in the
capacity-bound regime).} The tailored directions are nearly orthogonal to
$-\hat g$ (inner product $0.08$--$0.18$), so they are a genuine
alternative, yet they do not move OURS onto the CIFAR-10
stealthy-damage frontier. Across the $24$ cells
($4$ directions $\times 3$ $\sigma$ $\times 2$ $\kappa$), the uncapped
Q25 allocator inflicts strictly larger raw damage than OURS in $19$ cells;
OURS wins in only $5$, and \emph{all $5$ OURS wins occur at
$\kappa{=}2$ (loose budget)} where the stealth cap $\eta_{\max}$ is the
sole binding constraint (Table~\ref{tab:dirsens}, Fig.~\ref{fig:9}). The
binding constraint is therefore the \emph{per-client capacity} $C$, not
the attack direction. Conversely, the T3 monotonicity (frontier damage
$\uparrow$ with $\sigma$) is reproduced under every direction: the
best-attacker damage rises monotonically
$0.0187\!\to\!0.0210\!\to\!0.0274$ for $\sigma{=}\{0.5,1,2\}$, and three
of four directions are individually monotone. The fourth,
\texttt{shej\_sign}, shows a shallow dip at $\sigma{=}1$: per-seed drops
there are $\le 0.003$ pp under both OURS and Q25, and the effect size
versus $\sigma{=}2$ is $<0.005$ pp ($95\%$ CI width $<0.02$ pp), although
the paired Wilcoxon contrast is nominally significant at $n{=}5$ seeds in
all four (strategy,$\kappa$) contrasts ($p{=}0.002$--$0.049$); the
contrast versus $\sigma{=}0.5$ is significant in only $1$ of $4$
contrasts. We therefore read the dip as a noise-scale wiggle, not a
structural sub-$\epsilon$ feature.
The conclusion is robust to attack sophistication: the capacity-bound
regime is a property of the joint-defense geometry, not of a weak default
direction. (A faithful full re-implementation of Wang's surrogate-trained
LLRA and Hossain's RL privacy-level selector is left to future work; their
\emph{mechanisms} are covered by the directions above and the
$\eta_{\max}$ cap, which is Hossain's hide-within-the-noise rule.)

\subsubsection{R3: robustness to the aggregation rule}
\label{sec:r3}
Theory and measurement so far use coordinate-median aggregation and
median-based detection. We swap the aggregation rule inside the same
CIFAR-10 pipeline: coordinate median, Trimmed-Mean~\cite{yin2018}
($\beta{=}0.2$), and Krum~\cite{blanchard2017} ($f{=}6$), each with its
natural detection statistic (distance to the aggregate; Krum's
neighbor-discrepancy score), and re-probe the per-client stealthy
capacity $\eta_{\max}(\sigma)$ over
$\sigma_{\mathrm{rel}}\in\{0.5,\dots,16\}$ (no local training required;
Fig.~\ref{fig:10}). As a modern robust-aggregation baseline we add
FLAME~\cite{allouah2023flame} (iterative cosine filtering around the
coordinate median, accepted-set averaging, adaptive noise
$\delta_m{=}\lVert m\rVert/(f\sqrt d)$), scored with its natural
cosine-discrepancy statistic at the same $5\%$ benign FPR.

\emph{Finding (level changes; floor regime and $C/\Gamma$ boundary preserved).} All three aggregators
exhibit the floor of Proposition~1: the median
client's $\eta_{\max}$ is flat at $\eta_{\mathrm{f}}\approx1.15$--$1.23$
(variation $<7\%$) and the aggregate capacity in the floor regime is
$C=8.2$--$9.1$, i.e.\ the capacity-bound verdict $C/\Gamma<1$ is
aggregator-invariant. The crossover scale $\sigma_c$, however, is
aggregator-dependent: the median detector recovers
($\sigma_c{\approx}19$, $R^2{=}0.97$), Trimmed-Mean recovers more
slowly ($\sigma_c{\approx}27$, $R^2{=}0.74$), and Krum remains
floor-dominated throughout the probed window ($\sigma_c>16$); the lower
$R^2{=}0.74$ for Trimmed-Mean reflects a heavier-tailed score
distribution under $\beta{=}0.2$ trimming in the floor regime, which
widens the benign-threshold scatter and degrades the two-regime fit (a
fit-quality caveat, not a structural departure). The
mechanism is scale-equivariance: Krum's score statistic inflates
uniformly with $\sigma$, so the benign threshold rises with the noise
and the noise-dominated recovery is deferred beyond the probed window,
so the second regime is \emph{not observed} for Krum in this window. This
refines Section~7: swapping the aggregation
rule rescales when noise starts to help the attacker, but does not
change whether the floor regime and the $C/\Gamma$ boundary exist.

\emph{FLAME inverts the picture.} The cosine-discrepancy statistic is
insensitive to a perturbation along $-\hat g$ in this pipeline: benign
updates (norms $0.8$--$2.7$) are highly cosine-aligned, and even a
$\rho{=}12$ perturbation (norm $14.6$, five to eighteen times an entire
benign update) moves the score only from $0.854$ to $0.887$, below the
benign threshold $0.893$. The probed stealth capacities are therefore
right-censored at the probe's upper bound ($\eta_i{=}12$ in $17$ of
$36$ client-cells), the two-regime fit degrades ($R^2{=}0.31$,
$\eta_{\mathrm{f}}{=}6.05$), and the implied aggregate capacity
$C\ge53$ exceeds $\Gamma$ by an order of magnitude in \emph{both}
budget regimes. The aggregator-invariance established above is thus a
property of \emph{distance-to-aggregate} statistics: replacing distance
by cosine similarity does not rescale the floor but removes the floor's
detectability altogether, so the defender's $C/\Gamma$ verdict is not
aggregator-invariant across detection-statistic families. FLAME's
protection in this threat model reduces to its accepted-set averaging
and adaptive noise, which the end-to-end grid below shows leaves the
attack nearly unfiltered.

We then re-run the allocation comparison (OURS, Q25, W25) under each
of the three distance-to-aggregate aggregators across
$\sigma_{\mathrm{rel}}\in\{0.5,1,2\}$,
$\kappa\in\{0.25,2\}$, $5$ seeds ($270$ attack runs).

\emph{Finding (attack grid: regimes and stealth costs are
aggregator-invariant).} The budget-space structure of
Section~7.2 reproduces under every aggregator: the capacity ratio
is $C/\Gamma\approx3.7$--$4.1$ at $\kappa{=}0.25$ (deployable) versus
$0.47$--$0.52$ at $\kappa{=}2$ (capacity-bound). Whenever the cap
binds ($\kappa{=}2$), the uncapped Q25 allocator pays a uniformly
higher stealth price than OURS: its malicious-flag rate exceeds ours by
$8$--$28$ percentage points under all three aggregators, and its
per-round detection rate reaches $84$--$100\%$ versus $64$--$92\%$ for
OURS. Raw per-cell damage differences between allocators are small (max
$|\Delta|=1.70$ pp across the $18$ (aggregator,$\sigma$,$\kappa$) cells, median
$\approx 0.4$ pp), and paired Wilcoxon finds no cell significant at
$\alpha{=}0.05$ ($n{=}5$ seeds, all $p{\ge}0.062$). A pointwise bootstrap
additionally gives $4$ of $18$ nominal $95\%$ CIs that exclude $0$; these
intervals are \emph{unadjusted} for the $18$ comparisons, and after Holm
correction none remains significant, consistent with the Wilcoxon result. We
therefore report magnitudes rather than claim strict
seed-noise containment, and we do not claim per-cell damage-ordering invariance;
the aggregator-invariant
quantities are the regime boundary $C/\Gamma$ and the \emph{stealth
cost of ignoring the cap}. Krum ($f{=}6$) neutralizes the absolute
stakes (all drops within $\pm0.015$ of the no-attack control), yet the
same stealth-cost ordering persists: a stronger aggregator shrinks the
payoff matrix but does not change which allocator plays it better.

\emph{FLAME end-to-end.}
Re-running the same attack grid (OURS, Q25, W25;
$\sigma_{\mathrm{rel}}\in\{0.5,1,2\}$, $\kappa\in\{0.25,2\}$, $5$ seeds)
under FLAME confirms the probe: the cosine filter rarely fires
(round-level detection $0$--$40\%$, mean $19\%$, versus $36$--$100\%$
under the median rule in the same cells), FLAME's adaptive noise
$\delta_m=\lVert m\rVert/(f\sqrt d)$ is negligible relative to the
update scale, and the aggregation reduces to averaging over a highly
aligned accepted set. The end-to-end damage is accordingly positive in
\emph{every} cell ($0.013$--$0.128$): at $\kappa{=}0.25$ the accuracy
drop is $0.013$--$0.039$ versus $-0.002$--$0.024$ under the median
rule, and at $\kappa{=}2$ it is $0.094$--$0.128$ versus
$-0.001$--$0.040$, i.e.\ up to two orders of magnitude larger
damage in matched cells. FLAME is therefore a scope statement, not a rebuttal:
modern robust aggregators built on \emph{cosine} statistics do not
bound the stealthy capacity of per-update-DP-masked, median-aligned
attacks, whereas the distance-to-aggregate family does -- at the price
of the same two-regime exposure quantified above.

%======================================================================
\subsubsection{R4: robustness to an online adaptive attacker}
\label{sec:r4}

The stealth-constrained allocation in~Section~7.2 is derived against a
\emph{static} attacker who fixes the per-client budget $p$ once from an
offline estimate of $\{\eta_{\max,i}\}$ and never revises it. A natural
objection (raised in the Limitations) is that a \emph{reactive} attacker,
observing which malicious clients the defender flags each round, could
reallocate the total budget $\Gamma$ toward survivors and recover the
utility it would otherwise lose to detection. We close this gap with an
ablation on the CIFAR-10 pipeline.

\emph{Setup.} We keep the attack direction $-\hat g$ (identical to OURS)
and the same detection oracle (coordinate-median, $1.645\sigma$
threshold). The \textbf{static} attacker is exactly OURS: one
water-filling solve at $r{=}0$, budget held fixed for all $5$ rounds. The
\textbf{adaptive} attacker re-estimates, at every round $r{\ge}1$, the
per-client stealth cap: clients flagged in the previous round are retired
($\eta_i\!\leftarrow\!\eta_i/2$) and unflagged survivors are pushed toward
their offline cap ($\eta_i\!\leftarrow\!\min(\hat\eta_i,1.15\eta_i)$);
$\Gamma$ is then re-solved by water-filling over the surviving active set.
We sweep $\sigma_{\mathrm{rel}}\in\{0.5,1,2\}$ and
$\kappa\in\{0.25,2\}$ over $5$ seeds ($60$ attack runs per mode).

\emph{Finding (adaptive cannot beat the static optimum).} In all six
$(\sigma,\kappa)$ cells the adaptive attacker's accuracy drop is
\emph{not larger} than the static one; the signed gain
$g=\mathrm{drop}_{\mathrm{adapt}}-\mathrm{drop}_{\mathrm{static}}$ is
negative throughout, ranging from $-0.0028$ to $-0.0201$
(Fig.~\ref{fig:11}). Two regimes emerge. In the deployable regime
($\kappa{=}0.25$, few detections) static and adaptive are within seed
noise of each other ($|g|{\le}0.014$): with little to recover, adaptivity
is moot. In the capacity-bound regime ($\kappa{=}2$, many detections) the
adaptive attacker lowers its exposure, its per-round
detection rate falling from $96\%$--$68\%$ (static) to $64\%$--$48\%$,
but only by retiring flagged clients, which shrinks the active set and the
damage it can deliver: its malicious-flag rate drops to $17\%$--$22\%$ and
its accuracy drop turns non-positive (down to $-0.005$), i.e.\ the attack
ceases to degrade accuracy. The net effect is zero or negative gain.

\emph{Implication.} The static water-filling allocation is robust to an
online reactive attacker, not merely to a static proxy. The reason is
structural: the hard constraint is the capacity ceiling
$C=\sum_i\eta_{\max,i}^2$, and reallocating $\Gamma$ among survivors
cannot raise $C$; it can only trade detection exposure for potency, a
trade that, under a fixed $\Gamma$, never recovers the budget lost to
detection. This converts the Limitations caveat ``optimal against a static
attacker'' into a verified claim: the optimum is adaptively robust, and the
remaining open question is full online \emph{re-estimation} of
$\{\eta_{\max,i}\}$ (not reallocation), which we leave to future work.

%======================================================================

\paragraph{S2. Robustness: attack-direction sensitivity}
Table~\ref{tab:dirsens} and Fig.~\ref{fig:9} document the
direction-sensitivity analysis of \S\ref{sec:r2c}.
\begin{table}[!t]
\centering\small
\caption{Direction-sensitivity summary (CIFAR-10, $10$ seeds). Directions:
unit=agnostic, sign=static sign-flip, std=variance-tailored. Frontier
max drop = $\max$ over OURS/Q25 and $\kappa{\in}\{0.5,2\}$. Win counts
are over the $6$ $(\sigma,\kappa)$ cells per direction.}
\label{tab:dirsens}
\scriptsize\setlength{\tabcolsep}{3pt}
\begin{tabular}{lccc}
\toprule
Direction & Max drop @ $\sigma{=}.5/1/2$ & T3 mono? & Q25/OURS wins \\
\midrule
ours $-\hat g$            & 0.0187/0.0210/0.0274 & yes & 5 / 1 \\
Shej unit                 & 0.0036/0.0060/0.0131 & yes & 3 / 3 \\
Shej sign                 & 0.0040/0.0026/0.0117 & noisy & 5 / 1 \\
Shej std   & 0.0049/0.0078/0.0128 & yes & 6 / 0 \\
\emph{best attacker}& \emph{0.0187/0.0210/0.0274} & \emph{yes} & --- \\
\bottomrule
\end{tabular}
\end{table}
\begin{figure}[!t]
\centering
\includegraphics[width=0.95\linewidth]{figures/fig9_dirsens_cifar10.pdf}
\caption{Direction-sensitivity analysis. \textbf{A}: T3 monotonicity
reproduced under every attack direction; the best-attacker damage (black,
max over directions) rises monotonically with $\sigma$. \textbf{B}:
per-cell $(\text{Q25}-\text{OURS})$ raw damage; red = uncapped dominates
($19/24$ cells), blue = stealth-capped water-filling wins (all $5$ at
$\kappa{=}2$). The capacity-bound regime is attack-direction invariant.}
\label{fig:9}
\end{figure}
\paragraph{S3. Robustness: aggregator swap}
Fig.~\ref{fig:10} documents the aggregator-swap probe of \S\ref{sec:r3}.
\begin{figure}[!t]
\centering
\includegraphics[width=0.95\linewidth]{figures/fig10_aggswap_cifar10.pdf}
\caption{Aggregator-swap probe (CIFAR-10, R3-A). \textbf{A}: median
client's $\eta_{\max}(\sigma_{\mathrm{rel}})$ under coordinate median,
Trimmed-Mean ($\beta{=}0.2$), and Krum ($f{=}6$); all three are
floor-dominated with $\eta_{\mathrm{f}}\approx1.15$--$1.23$.
\textbf{B}: two-regime fits; the floor level is aggregator-invariant
while the crossover $\sigma_c$ is not (Krum's scale-equivariant score
defers recovery beyond the probed window).}
\label{fig:10}
\end{figure}

\paragraph{S4. Robustness: online adaptive attacker}
Fig.~\ref{fig:11} documents the adaptive-attacker ablation of \S\ref{sec:r4}.
\begin{figure}[!t]
\centering
\includegraphics[width=\linewidth]{figures/fig11_adaptive.pdf}
\caption{Online adaptive-attacker ablation (CIFAR-10, R4). Accuracy drop
vs no-attack for the static (OURS) and adaptive attackers across DP levels
and defense budgets. The signed gain $\Delta$ (adaptive $-$ static) is
negative in every cell: a reactive attacker that reallocates $\Gamma$
toward survivors cannot exceed the static optimum, because the capacity
ceiling $C$ binds.}
\label{fig:11}
\end{figure}

\subsubsection{R5: cross-dataset validation figure (MNIST)}
The cross-dataset validation discussed in the R5 subsection of the
main text is reproduced here.

\begin{figure}[!t]
\centering
\includegraphics[width=\linewidth]{figures/fig12_crossdataset.pdf}
\caption{Cross-dataset validation (R5). \textbf{A}: the two-regime
$\eta_{\max}(\sigma)$ law (plateau at low $\sigma$, refinement at high
$\sigma$) holds on both CIFAR-10 and MNIST; the naive
$\eta\propto\sigma$ reference (dashed) is violated on both. \textbf{B}:
under the $\kappa{=}2$ setting (capacity-bound for CIFAR-10) the accuracy
drop is
dataset-dependent (negligible on MNIST with $97.7\%$ clean accuracy,
$0.2$--$0.8\%$ on CIFAR-10), yet the structural stealth-cost ordering
of the allocators is preserved on both.}
\label{fig:12}
\end{figure}


\clearpage
\begin{thebibliography}{99}
\bibitem{shejwalkar2021turning}
V. Shejwalkar and A. Houmansadr, ``Manipulating the Byzantine: Optimizing
model poisoning attacks and defenses for federated learning,'' in
\emph{Proc. NDSS}, 2021.
\bibitem{hossain2025mpelm}
M.~T. Hossain, S. Badsha, H. La, S. Islam, and I. Khalil, ``Exploiting
Gaussian noise variance for dynamic differential poisoning in federated
learning,'' \emph{IEEE Trans.\ Artificial Intelligence}, vol.~6, no.~11,
pp. 2922--2939, Nov. 2025, doi:10.1109/TAI.2025.3540030.
\bibitem{wang2025ldpfl}
Z. Wang, W. Tong, T. Han, H. Chen, T. Zhang, Y. Mao, and S. Zhong, ``On
evaluating the poisoning robustness of federated learning under local
differential privacy,'' arXiv:2509.05265, 2025. The code repository
announced in the paper (\url{https://github.com/ZiJW/LDPFL-Attack}) is
no longer publicly accessible as of Oct.\ 2026.
\bibitem{qiu2025budget}
J. Qiu, Y. Zhao, and J. Li, ``Optimal perturbation budget allocation for
data poisoning in offline reinforcement learning,'' arXiv:2512.08485,
Dec. 2025.
\bibitem{mcmahan2017}
H.~B. McMahan \etal, ``Communication-efficient learning of deep networks
from decentralized data,'' in \emph{Proc. AISTATS}, 2017.
\bibitem{blanchard2017}
P. Blanchard \etal, ``Machine learning with adversaries: Byzantine
tolerant gradient descent,'' in \emph{Proc. NeurIPS}, 2017.
\bibitem{allouah2023flame}
R. Allouah, A. Farhadkhani, R. Guerraoui, N. Gupta, R. Pinot, and
J. Stephan, ``Fixing by mixing: A recipe for optimal byzantine ML under
heterogeneity,'' in \emph{Proc. ICML}, 2023.
\bibitem{yin2018}
D. Yin \etal, ``Byzantine-robust distributed learning: Towards optimal
statistical rates,'' in \emph{Proc. ICML}, 2018.
\bibitem{fang2020local}
M. Fang, X. Cao, J. Jia, and N.~Z. Gong, ``Local model poisoning attacks
to Byzantine-robust federated learning,'' in \emph{Proc. USENIX Security
Symposium}, 2020.
\bibitem{abadi2016}
M. Abadi \etal, ``Deep learning with differential privacy,'' in
\emph{Proc. CCS}, 2016.
\bibitem{kairouz2021advances}
P. Kairouz \etal, ``Advances and open problems in federated learning,''
\emph{Foundations and Trends in Machine Learning}, vol.~14, no.~1--2,
pp. 1--210, 2021.
\bibitem{li2020federated}
T. Li, A.~K. Sahu, A. Talwalkar, and V. Smith, ``Federated learning:
Challenges, methods, and future directions,'' \emph{IEEE Signal
Processing Magazine}, vol.~37, no.~3, pp. 50--60, 2020.
\bibitem{mothukuri2021survey}
V. Mothukuri, R.~M. Parre\~no, S. Gharehbaghi, H.~M. Huang, and Y. Si,
``A survey on security and privacy of federated learning,'' \emph{Future
Generation Computer Systems}, vol.~115, pp. 619--640, 2021.
\bibitem{biggio2012poisoning}
B. Biggio, B. Nelson, and P. Laskov, ``Poisoning attacks against
support vector machines,'' in \emph{Proc. ICML}, 2012.
\bibitem{meizhu2015teaching}
S. Mei and X. Zhu, ``Using machine teaching to identify optimal
training-set attacks on machine learners,'' in \emph{Proc. AAAI}, 2015.
\bibitem{biggio2018wild}
B. Biggio and F. Roli, ``Wild patterns: Ten years after the rise of
adversarial machine learning,'' \emph{Pattern Recognition}, vol.~84,
pp. 317--331, 2018.
\bibitem{bhagoji2019analyzing}
A.~N. Bhagoji, S. Chakraborty, P. Mittal, and S. Calo, ``Analyzing
federated learning through an adversarial lens,'' in \emph{Proc. ICML},
2019.
\bibitem{bagdasaryan2020backdoor}
E. Bagdasaryan, A. Veit, Y. Hua, D. Estrin, and V. Shmatikov, ``How to
backdoor federated learning,'' in \emph{Proc. AISTATS}, 2020.
\bibitem{xie2020dba}
C. Xie, K. Huang, P.-Y. Chen, and B. Li, ``DBA: Distributed backdoor
attacks against federated learning,'' in \emph{Proc. ICLR}, 2020.
\bibitem{sun2019can}
Z. Sun, P. Kairouz, and A.~T. Suresh, ``Can you really backdoor
federated learning?'' arXiv:1911.07963, 2019.
\bibitem{mhamdi2018hidden}
E.~M. El Mhamdi, R. Guerraoui, and S. Rouault, ``The hidden
vulnerability of distributed learning in Byzantium,'' in
\emph{Proc. ICML}, 2018.
\bibitem{bernstein2018signsgd}
J. Bernstein, Y.-X. Wang, K. Azizzadenesheli, and A. Anandkumar,
``signSGD: With majority vote is fast and error tolerant,'' in
\emph{Proc. ICLR}, 2018.
\bibitem{cao2021fltrust}
X. Cao, M. Fang, J. Liu, and N.~Z. Gong, ``FLTrust: Byzantine-robust
federated learning via trust bootstrapping,'' in \emph{Proc. NDSS},
2021.
\bibitem{nguyen2022flame}
T.~D. Nguyen \etal, ``FLAME: Taming backdoors in federated learning,''
in \emph{Proc. NeurIPS}, 2022.
\bibitem{xu2024minvar}
H. Xu and T. Shu, ``Defending against model poisoning attack in
federated learning: A variance-minimization approach,'' \emph{Journal
of Information Security and Applications}, vol.~82, art.~103744, 2024.
\bibitem{xu2024agnostic}
H. Xu and T. Shu, ``Attack-model-agnostic defense against model
poisonings in distributed learning,'' \emph{Journal of Information
Security and Applications}, vol.~82, 2024.
\bibitem{luong2024fedlsae}
T.~D. Luong, V.~M. Tien, N.~H. Quyen, D.~T.~T. Hien, P.~T. Duy, and
V.-H. Pham, ``Fed-LSAE: Thwarting poisoning attacks against federated
cyber threat detection system via autoencoder-based latent space
inspection,'' \emph{Journal of Information Security and Applications},
vol.~87, art.~103916, 2024.
\bibitem{mcmahan2018dp}
H.~B. McMahan, D. Ramage, K. Talwar, and L. Zhang, ``Learning
differentially private recurrent language models,'' in \emph{Proc.
ICLR}, 2018.
\bibitem{geyer2017dp}
R.~C. Geyer, T. Klein, and M. Nabi, ``Differentially private
multi-client federated learning,'' arXiv:1712.07557, 2017.
\bibitem{dwork2006calibrating}
C. Dwork, F. McSherry, K. Nissim, and A. Smith, ``Calibrating noise to
sensitivity in private data analysis,'' in \emph{Proc. TCC}, 2006.
\bibitem{dwork2014algorithmic}
C. Dwork and A. Roth, ``The algorithmic foundations of differential
privacy,'' \emph{Foundations and Trends in Theoretical Computer
Science}, vol.~9, no.~3--4, pp. 211--407, 2014.
\bibitem{bonawitz2017practical}
K. Bonawitz \etal, ``Practical secure aggregation for privacy-preserving
machine learning,'' in \emph{Proc. CCS}, 2017.
\bibitem{song2019privacy}
L. Song, R. Shokri, and P. Mittal, ``Privacy risks of securing machine
learning models against adversarial examples,'' in \emph{Proc. IEEE
Symposium on Security and Privacy}, 2019.
\bibitem{boyd2004convex}
S. Boyd and L. Vandenberghe, \emph{Convex Optimization}. Cambridge
University Press, 2004.
\bibitem{huber1964robust}
P.~J. Huber, ``Robust estimation of a location parameter,'' \emph{Annals
of Mathematical Statistics}, vol.~35, no.~1, pp. 73--101, 1964.
\bibitem{tambe2011security}
M. Tambe, \emph{Security and Game Theory: Algorithms, Deployed Systems,
Lessons Learned}. Cambridge University Press, 2011.
\bibitem{zhou2024roppfl}
H. Zhou, Y. Zheng, and X. Jia, ``Towards robust and privacy-preserving
federated learning in edge computing,'' \emph{Computer Networks},
vol.~243, art.~110321, 2024, doi:10.1016/j.comnet.2024.110321.
\bibitem{zhao2025fedmp}
K. Zhao, L. Wang, F. Yu, B. Zeng, and Z. Pang, ``FedMP: A multi-pronged
defense algorithm against Byzantine poisoning attacks in federated
learning,'' \emph{Computer Networks}, vol.~257, art.~110990, 2025,
doi:10.1016/j.comnet.2024.110990.
\bibitem{liu2025antidotefl}
Y. Liu, H. Zhang, M. Wang, Q. Xie, and Z. Sun, ``AntidoteFL: Enhancing
defense against poisoning attacks in federated learning,'' \emph{Computer
Networks}, vol.~269, art.~111427, 2025, doi:10.1016/j.comnet.2025.111427.
\end{thebibliography}

\end{document}
